\chapter*{Závěr}
\addcontentsline{toc}{chapter}{Závěr}

Předkládaná práce si položila otázku, jak algoritmicky zprostředkovaná digitální infrastruktura mění podmínky možnosti veřejné racionality a jaký druh politického problému tato změna představuje. 
Odpověď, k níž jsem dospěla, lze shrnout do čtyř tezí.


Za prvé: ztotožnění myšlení s výpočtem je pojmovým předpokladem algoritmické moci, nikoli jejím důsledkem. Rekonstrukce provedená v první kapitole ukázala, že cesta od Lullovy Ars Magna k současným jazykovým modelům je dějinami postupného vylučování: z pojmu myšlení byl nejprve vyloučen mravní úsudek, poté intencionalita, poté vtělenost a nakonec samotné porozumění. 
Politický význam tohoto vyloučení nespočívá v tom, že by stroje „nemyslely“ (spor o to nechávám otevřený) nýbrž v tom, že takto zúžený pojem rozumu učinil myslitelným delegování rozhodnutí na entitu, jež nesplňuje podmínky odpovědného subjektu.
 Depolitizace tedy není vedlejším účinkem technického pokroku; je zabudována do pojmové architektury, jež jej umožnila.


 Za druhé: platformový kapitalismus systematicky podkopává materiální předpoklady veřejné racionality, avšak rozsah této škody je sporný a nemůže nést tíhu politického argumentu. 
Druhá kapitola ukázala mechanismy: komodifikaci pozornosti, ohrazování zkušenosti do podoby dat, reálné podřízení komunikace logice kapitálu, erozi sdíleného epistemického prostoru.
 Zároveň jsem se však nemohla vyhnout tomu, že řada empirických tvrzení, o něž se kritická literatura opírá — kognitivní atrofie, epistemické bubliny, účinnost mikrotargetingu — je v odborné diskusi zpochybňována.
 Poctivá odpověď zní: nevíme s jistotou, jak velká ta škoda je. 
 Kritika, jež na těchto tvrzeních staví, je proto vratká.


 Za třetí, a to je hlavní teze práce: politické bezpráví algoritmické moci nespočívá ve škodě, nýbrž v dominanci.
Ať už je měřitelný kognitivní dopad platforem jakýkoli, trvá skutečnost, že nevolené soukromé subjekty disponují svévolnou, neprůhlednou a demokraticky nekontrolovanou mocí nad infrastrukturou, v níž se veřejná racionalita jedině může uskutečňovat: rozhodují, co je viditelné, komu jsou přiznány zdroje a příležitosti, a jaké epistemické prostředí občan obývá.
 V republikánském pojetí svobody jako nepřítomnosti dominance je toto bezpráví plné a nezávislé na tom, zda je moc uplatňována laskavě.
 Právě tento posun umožňuje vysvětlit, proč individualisticky pojatá ochrana práv — souhlas, přístup k datům, právo na vysvětlení — strukturálně selhává: dominance je vztah, nikoli událost, a nelze ji odstranit ani sebedokonalejším informovaným souhlasem jednotlivce.


 Za čtvrté: mezera odpovědnosti není technickým defektem, jenž by šlo technicky zaplnit.
 Rozostření odpovědnosti mezi vývojáře, provozovatele, zadavatele a regulátora reprodukuje logiku, kterou Arendtová popsala jako banalitu zla — nikoli jako historickou charakteristiku Eichmanna, nýbrž jako ideální typ byrokratické atomizace úkolu, v níž každý plní svou dílčí povinnost a souhrnný výsledek nemá původce.
 Antropomorfizace algoritmických systémů tuto mezeru nezaplňuje, nýbrž maskuje: připisuje subjektivitu artefaktu, jenž nesplňuje její podmínky, a osvobozuje tak skutečné subjekty od reflexe vlastního jednání. Odpovědnost však nelze delegovat na to, co nemá tvář.
\section*{Co z toho plyne}
Z takto vymezené diagnózy plynou tři odpovědi, jež se navzájem nevylučují, nýbrž doplňují. Ochrana (GDPR, AI Act) je nezbytná, avšak reaktivní, individualistická a asymetrická — a jak ukazuje osud tzv. 
Digital Omnibus, jímž bylo v roce 2026 uplatňování klíčových ustanovení AI Actu odloženo dříve, než vůbec vstoupilo v účinnost, je také vystavena systematickému tlaku regulovaných.
 Rozdělení moci — antimonopolní politika, povinný audit, oddělení infrastruktury od služeb — odpovídá povaze problému lépe, neboť míří na strukturu, nikoli na jednotlivý účinek.
 A konečně zmocnění: budování alternativní, federativní a společně vlastněné digitální infrastruktury, jež logiku dominance nereprodukuje.
Právě zde je však na místě střízlivost, k níž mne analýza dovedla. 
Existující alternativy — Fediverse, Wikipedia, otevřené jazykové modely, participativní platformy typu Decidim — dokládají, že jiná institucionální forma je možná. Nedokládají však, že je dostatečná.
 Migrace uživatelů na Mastodon po roce 2022 se z hlediska udržení uživatelů ukázala jako podstatně méně úspěšná, než naznačovala její první vlna; 
 participace na deliberativních platformách zůstává nízká; otevřené modely jsou závislé na výpočetní kapacitě, kterou nekontrolují. 
 Alternativní infrastruktura je nutnou, nikoli postačující podmínkou. 
 To není důvod k rezignaci, nýbrž k přesnějšímu zadání: otázka nezní, zda lze postavit lepší platformu, nýbrž jaké institucionální a majetkové uspořádání by zabránilo tomu, aby se moc znovu zkoncentrovala.
\section*{Meze práce a další výzkum}
Práce je prací pojmovou; empirické otázky, jichž se dotýká, nemůže rozhodnout.
 Přijímá habermasovský normativní rámec, ačkoli ten je sám předmětem sporu — agonistická kritika (Mouffe) by namítla, že ideál racionálního konsensu depolitizuje stejně účinně jako technokracie, proti níž je namířen; 
 tuto námitku jsem mohla pouze naznačit, nikoli vyřešit. 
Neanalyzuje mimoevropské konfigurace algoritmické moci ani otázku výpočetní suverenity, jež se pravděpodobně stane ústředním politickým problémem následujícího desetiletí.
Přesto se domnívám, že hlavní výsledek obstojí. Otázka, kterou algoritmická moc klade demokracii, není otázkou psychologickou ani technickou.
 Je to nejstarší otázka politické filozofie, položená v novém materiálu: kdo drží moc, komu se z ní zodpovídá a jak ji lze zpochybnit. 
 Kantovo Sapere aude nebylo adresováno géniovi schopnému samostatné myšlenky za jakýchkoli podmínek, nýbrž obyčejnému občanu, jehož schopnost k vlastnímu úsudku je strukturálně závislá na podmínkách, za nichž se formuje.
  Tyto podmínky dnes nevytváří veřejnost. Vytváří je několik soukromých korporací — a činí tak bez pověření, bez odpovědnosti a bez možnosti odvolání. 
  Obnovit veřejnou racionalitu proto neznamená naučit občany lépe zacházet s technologií. 
  Znamená to vrátit rozhodování o infrastruktuře veřejného rozumu do rukou těch, jichž se týká.
